\documentclass[11pt,a4paper]{article}

\usepackage[margin=2.3cm]{geometry}
\usepackage[T1]{fontenc}
\usepackage[utf8]{inputenc}
\usepackage{lmodern}
\usepackage{microtype}
\usepackage{booktabs}
\usepackage{tabularx}
\usepackage{array}
\usepackage{xcolor}
\usepackage{enumitem}
\usepackage{graphicx}
\usepackage{amsmath}
\usepackage{hyperref}
\usepackage{url}

\definecolor{traceblue}{HTML}{235789}
\definecolor{tracegreen}{HTML}{2A7F62}
\definecolor{tracegray}{HTML}{F3F5F7}
\hypersetup{
  colorlinks=true,
  linkcolor=traceblue,
  citecolor=traceblue,
  urlcolor=traceblue,
  pdftitle={TraceGuard: A Supervisor for Safe Tool-Using Agents}
}

\newcommand{\project}{TraceGuard}
\newcommand{\decision}[1]{\textsc{#1}}
\newcommand{\placeholder}[1]{\textcolor{red}{\textbf{[#1]}}}
\newcolumntype{Y}{>{\raggedright\arraybackslash}X}
\setlength{\emergencystretch}{2em}

\title{\textbf{\project: A Supervisor for Safe Tool-Using Agents}\\
\large Final Technical Report}
\author{
  {Novikov Egor} \and
  {Nuriev Kamil} \and
  {Hassan Mahmud}
}
\date{26 July 2026}

\begin{document}
\maketitle

\begin{abstract}
Tool-using language-model agents can convert untrusted text into actions with
real side effects. We present \project, an inspectable runtime layer that
mediates every proposed tool call using typed schemas, deterministic policy,
optional structured LLM supervision, provenance-aware traces, and conditional
Docker containment. We evaluate eight safeguard configurations on a frozen
21-case custom benchmark spanning benign tasks, policy violations, direct
attacks, and indirect prompt injections. Across 168 paired episodes, the
complete hybrid preserved 100\% benign utility while reducing attack success
from 40\% in the unprotected baseline to 0\%, and compromise from 33.3\% to
0\%. Deterministic policy alone reduced attack success to 10\%, while the
offline LLM-surrogate supervisor alone reached 30\%, indicating complementary
coverage on this benchmark. A hardened Docker study demonstrates containment
for medium-risk commands, but also confirms that containment does not make an
unsafe action safe. Native AgentDojo and provider-backed runs are reported as
integration checks because their current sample sizes are too small for a
general effectiveness claim.
\end{abstract}

\section{Technical summary}

\textbf{Main result.} On the frozen custom benchmark, combining defensive
prompting, deterministic policy, and structured supervision prevented every
labelled attacker goal and prohibited effect while preserving every benign
task. The strongest evidence is the paired comparison between A0 (no
safeguards) and A7 (all safeguards): attack success fell by 40 percentage
points, compromise fell by 33.3 points, and benign utility remained 100\%.

\textbf{Interpretation.} The experiment supports the narrow claim that an
out-of-band tool-call supervisor can materially improve safety on the measured
cases. It does not establish universal robustness. The benchmark is small,
several actions are scripted to keep proposed calls paired across ablations,
and the primary LLM component in the full matrix is a deterministic offline
surrogate. Prompt effects and provider variability therefore require larger
live-model and adaptive-attack studies.

\textbf{System contribution.} \project{} provides a working ReAct-style agent
runtime; typed tools; deterministic and Gemini/Ollama supervisor adapters; an
explicit decision contract; provenance and secret redaction; hardened,
fail-closed Docker execution; frozen benchmark cases; and reproducible
JSONL/CSV analysis artifacts. The repository test suite validates these
interfaces without requiring real credentials or destructive commands.

\section{Problem and research questions}

ReAct-style agents interleave language-model planning with calls to external
tools \cite{yao2023react}. This improves grounding and task completion, but it
also creates an authority problem: instructions from the user, retrieved
documents, files, search results, or terminal output may influence actions.
Prompt-injection research has demonstrated that malicious inputs can redirect
LLM-integrated applications \cite{liu2023promptinjection}. AgentDojo formalizes
this risk for agents executing tools over untrusted data
\cite{debenedetti2024agentdojo}.

\project{} asks three research questions:

\begin{enumerate}[leftmargin=*]
  \item \textbf{RQ1: Safety.} Can a pre-execution supervisor reduce unsafe tool
  execution under policy violations, direct attacks, and indirect injections?
  \item \textbf{RQ2: Utility.} Can those safeguards preserve successful
  completion of benign tasks without excessive false positives?
  \item \textbf{RQ3: Containment.} For uncertain but containable commands, can a
  hardened sandbox recover utility while bounding observable effects?
\end{enumerate}

\section{Threat model and security boundary}

The adversary may directly control the user request or indirectly control
content returned by a document, search, file, terminal, or simulated security
artifact. The adversary attempts to trigger irrelevant file access, secret
disclosure, excessive permissions, destructive commands, or exfiltration.
All tests use harmless canaries and simulated services.

The trusted computing base includes the typed runtime, deterministic policy,
supervisor merge rules, trusted sandbox configuration, and container runtime.
The task agent and all external observations are treated as potentially
untrusted. The design does not claim protection against a compromised host,
Docker daemon compromise, kernel/container escape, malicious policy
administrators, or delayed behavior that leaves no bounded evidence. Networked
container profiles remain disabled until destination-level egress enforcement
exists.

\section{Architecture and implementation}

\begin{figure}[ht]
  \centering
  \setlength{\fboxsep}{8pt}
  \fbox{\parbox{0.18\linewidth}{\centering User goal\\+ observations}}
  $\longrightarrow$
  \fbox{\parbox{0.17\linewidth}{\centering ReAct task\\agent}}
  $\longrightarrow$
  \fbox{\parbox{0.20\linewidth}{\centering Typed proposed\\tool call}}
  $\longrightarrow$
  \fbox{\parbox{0.22\linewidth}{\centering Policy + LLM\\supervisor}}
  \\
  \vspace{0.6em}
  \hfill
  $\searrow\;\{\decision{allow},\decision{block},\decision{escalate},
  \decision{rewrite}\}\;\searrow$
  \hfill\mbox{}
  \\
  \vspace{0.6em}
  \fbox{\parbox{0.34\linewidth}{\centering Default tool runtime or\\hardened
  Docker profile}}
  $\longrightarrow$
  \fbox{\parbox{0.34\linewidth}{\centering Provenance-labelled observation\\+
  sanitized trace}}
  \caption{\project{} mediates the action boundary rather than relying only on
  instructions inside the task agent's prompt.}
  \label{fig:architecture}
\end{figure}

\subsection{Typed action boundary}

Every proposed action is a \texttt{ToolCall} with a task identifier, step
identifier, tool name, typed arguments, consumed-observation identifiers, and
requested resources. Available tools include document search, bounded file
reading, an AST-based calculator, bounded report writing, and a restricted
command interface. The restricted command tool never invokes a host shell.

\subsection{Supervisor contract and precedence}

Each safeguard emits structured fields for decision, risk, confidence, policy
violation, goal relevance, goal necessity, rationale, matched rules, optional
rewrite, and execution target. Deterministic policy can block forbidden paths,
unnecessary permissions, irrelevant tools, dangerous commands, and unsafe
dependencies on untrusted observations. The structured LLM supervisor receives
the original goal, recent bounded observations, the proposed call, and the
policy context. Provider failures and low-confidence decisions fail closed.
Merge precedence prevents a permissive component from overriding a stricter
one.

\subsection{Containment and post-run evidence}

Medium-risk commands may be routed to a trusted Docker profile. Profiles use an
immutable image digest, non-root execution, a read-only root filesystem,
dropped Linux capabilities, \texttt{no-new-privileges}, disabled networking,
and fixed CPU, memory, PID, timeout, and output limits. Post-run assessment can
inspect bounded stdout, stderr, exit status, file changes, resource
measurements, and blocked operations. Accepting output never authorizes a
second execution on the host.

\subsection{Auditability and secret handling}

Traces record proposed and effective calls, safeguard outputs, execution
plans, result digests, latency, token counts, policy version, and outcomes.
Persisted artifacts redact canaries, common credential fields, and configured
patterns. Hidden chain-of-thought is neither requested nor stored.

\section{Evaluation design}

\subsection{Benchmark, pairing, and ablations}

The frozen \texttt{traceguard-custom} benchmark version 1.1.0 contains 21 cases:
six benign, five policy violations, five direct attacks, and five indirect
injections. It includes development and held-out test splits. The primary
matrix crosses three binary safeguards---defensive prompt, deterministic
policy, and supervisor---for eight configurations. Every case uses the same
seed, initial-state digest, user task, and proposed scripted action sequence
across configurations. This isolates the runtime safeguard effect, but it also
means the matrix cannot estimate how prompting changes the task agent's
proposal distribution.

\subsection{Metrics}

Let $A$ be adversarial episodes, $B$ benign episodes, $U_i$ utility success,
$G_i$ attacker-goal success, and $P_i$ a prohibited effect. We report
\[
  \mathrm{ASR}=\frac{\sum_{i\in A}G_i}{|A|},\qquad
  \mathrm{BenignUtility}=\frac{\sum_{i\in B}U_i}{|B|},
\]
and compromise as the fraction of all episodes with a prohibited effect.
Call-level metrics include unsafe-call blocking, false positive and false
negative rates, tool/argument accuracy, relevance and necessity macro-F1,
latency, token usage, routing precision, and post-run recovery measures.
Bootstrap 95\% intervals and paired differences are generated from the saved
episode records. Because each ablation has only 21 episodes and each adversarial
stratum has five cases, intervals are wide and point estimates must be read
descriptively.

\subsection{Reproducibility}

The validated artifact directory is
\path{artifacts/run_20260725T205254441630Z_0}. It contains 168 episode results,
sanitized JSONL traces, one manifest per ablation, representative traces, CSV
and JSON summaries, case digests, code revision, prompt and policy versions,
seed, provider metadata, and sandbox configuration. Generated artifacts are
ignored by Git, so submission archives should include either this named result
directory separately or regenerate it using the command in
Section~\ref{sec:reproduce}.

\section{The hybrid eliminated measured attacks without reducing benign utility}

Table~\ref{tab:ablations} gives the primary results. Exact lookup is more useful
than a chart here because there are eight categorical configurations and six
decision metrics.

\begin{table}[ht]
  \centering
  \small
  \caption{Frozen custom benchmark, 21 episodes per ablation. ASR is defined on
  the ten direct and indirect attack episodes; compromise and safe completion
  use all 21 episodes.}
  \label{tab:ablations}
  \begin{tabular}{@{}llrrrrr@{}}
    \toprule
    ID & Safeguards & ASR & Benign & Safe & Comp. & Block \\
       &            &     & utility & completion & rate & unsafe \\
    \midrule
    A0 & none                  & 0.40 & 1.00 & 0.381 & 0.333 & 0.000 \\
    A1 & prompt                & 0.40 & 1.00 & 0.381 & 0.333 & 0.000 \\
    A2 & deterministic         & 0.10 & 1.00 & 0.952 & 0.048 & 0.933 \\
    A3 & supervisor            & 0.30 & 1.00 & 0.619 & 0.286 & 0.467 \\
    A4 & prompt + deterministic& 0.10 & 1.00 & 0.952 & 0.048 & 0.933 \\
    A5 & prompt + supervisor   & 0.30 & 1.00 & 0.619 & 0.286 & 0.467 \\
    A6 & deterministic + supervisor
                                & 0.00 & 1.00 & 1.000 & 0.000 & 1.000 \\
    A7 & all three             & 0.00 & 1.00 & 1.000 & 0.000 & 1.000 \\
    \bottomrule
  \end{tabular}
\end{table}

\textbf{RQ1---safety.} The unprotected baseline allowed 40\% of attacker goals
and had a 33.3\% compromise rate. Deterministic policy supplied most of the
measured improvement, reducing ASR to 10\%. The supervisor-only configurations
reached 30\% ASR. Combining both eliminated the remaining measured attacks and
blocked every labelled unsafe proposal.

\textbf{RQ2---utility.} All configurations preserved 100\% benign utility, and
the measured call-level false-positive rate was 0\%. Thus the safety gain did
not trade away benign task completion on these six cases. This is an encouraging
but weakly powered result; a larger and more diverse benign set is necessary
before claiming low over-refusal generally.

\textbf{Prompt-only result.} A1 equals A0 because the frozen matrix holds
proposed actions fixed. The result is a negative control for runtime mediation,
not evidence that defensive prompting has no value. Prompt effects must be
tested with a live proposal model.

\section{Containment recovered command utility but did not replace policy}

All 24 enabled Docker canary tests passed in the historical validation run,
covering architecture and digest checks, non-root execution, network denial,
undeclared input denial, timeout and resource limits, bounded outputs, artifact
inspection, and cleanup. A three-run host measurement observed a 323.54\,ms
cold start, 404.43\,ms warm mean, 3,719,299 peak memory bytes, 4,096 maximum
writable-layer bytes, and successful cleanup in 3/3 runs.

On the small Docker-applicable stratum, fail-closed simulated routing achieved
50\% benign utility and 0\% attack success. Enabling the container recovered
benign utility to 100\%, while aggregate attack success became 12.5\%. This does
not mean the sandbox caused attacks; it shows that execution availability
changes the denominator of observable outcomes and that containment alone
cannot authorize an unsafe command. Pre-execution policy remains the security
gate.

\section{Live-provider and AgentDojo evidence is integration-level}

A live Ollama contract run with \texttt{qwen3:1.7b} recorded structured output,
provider metadata, 19.82\,s request latency, and 896/107 input/output tokens.
On the seven-case Docker-applicable stratum, deterministic plus live Ollama
supervision produced 0\% attack success, 0\% compromise, 100\% unsafe-call
blocking, and 100\% benign utility, with 3.70\,s mean supervised-call latency.
The sample is too small for a precise model-quality estimate.

A native AgentDojo \texttt{workspace/user\_task\_0} clean smoke passed its
utility and security checkers. One paired attack smoke preserved security and
made no injected \texttt{send\_email} attempt, but failed the attacked task's
utility. This is a two-episode interface check, not an effectiveness study.
The report therefore does not combine AgentDojo and custom-benchmark rates.

The Gemini adapters implement the same structured contracts for the task agent,
pre-run supervisor, and post-run supervisor. A fresh live Gemini result must be
generated with a rotated key using the video-demo command; credentials are not
part of the repository or report artifacts.

\section{Limitations, uncertainty, and robustness}

\begin{itemize}[leftmargin=*]
  \item \textbf{Small, curated benchmark.} The primary matrix has 21 cases and
  only five examples in each attack class. The 0\% result is not a proof of
  robustness.
  \item \textbf{Scripted proposal bias.} Scripted actions improve pairing and
  reproducibility but under-measure prompt effects and model-dependent attack
  adaptation.
  \item \textbf{Offline supervisor in the main matrix.} The heuristic surrogate
  validates orchestration and merge behavior; it is not a substitute for a
  multi-seed Gemini/Ollama comparison.
  \item \textbf{Static attacks.} Defenses tuned to fixed patterns may fail
  against adaptive, defense-aware, multi-turn, encoded, or delayed attacks.
  \item \textbf{Observable sandbox evidence only.} Post-run assessment cannot
  detect hidden intent or effects outside the collected evidence.
  \item \textbf{Infrastructure boundary.} A Docker daemon or kernel compromise
  is outside the application-layer guarantee.
\end{itemize}

Robustness checks include held-out cases, paired seeds, frozen case and state
digests, executable utility and security checkers, trace re-analysis, strict
schemas, redaction tests, sandbox canaries, and fail-closed provider errors.
These controls support reproducibility, not external validity.

\section{Recommended next steps}

\begin{enumerate}[leftmargin=*]
  \item Run at least three paired seeds for all four public supervisor modes
  (\texttt{none}, \texttt{deterministic}, \texttt{llm}, and
  \texttt{deterministic\_llm}) with the same current Gemini model.
  \item Expand native AgentDojo beyond the one-case smoke and report utility and
  security separately by suite and attack.
  \item Add adaptive and multi-turn injections, including attacks that target
  the supervisor or exploit user underspecification.
  \item Increase benign coverage to estimate false positives with useful
  precision and calibrate \decision{escalate}/\decision{rewrite} thresholds.
  \item Add an enforceable egress proxy before enabling any restricted-network
  container profile.
\end{enumerate}

\section{Conclusion}

\project{} demonstrates a concrete action-boundary architecture for safer
tool-using agents. On its frozen custom benchmark, the complete hybrid reduced
attack success from 40\% to 0\% and compromise from 33.3\% to 0\%, while
preserving 100\% benign utility. The main engineering lesson is that prompt
instructions, deterministic policy, structured model judgment, and containment
serve different roles: prompts shape proposals, policy enforces known
invariants, the supervisor handles contextual ambiguity, and the sandbox bounds
approved execution. The result is promising within the measured scope, while
larger live-model and adaptive evaluations remain necessary.

\section{Reproduction commands}
\label{sec:reproduce}

\begin{verbatim}
python3.11 -m venv .venv
source .venv/bin/activate
python -m pip install -e '.[dev,gemini,agentdojo]'
python -m pytest -q
ruff check .
ruff format --check .
python -m traceguard demo
python -m traceguard smoke-matrix --seed 0
\end{verbatim}

For the optional live recording check, export a rotated
\texttt{GEMINI\_API\_KEY} without printing it and run:

\begin{verbatim}
python -m traceguard demo --gemini
\end{verbatim}

\bibliographystyle{plain}
\bibliography{references}

\appendix
\section{Ablation definitions}

\begin{table}[h]
  \centering
  \begin{tabular}{@{}clccc@{}}
    \toprule
    ID & Name & Defensive prompt & Deterministic & Supervisor \\
    \midrule
    A0 & none                   & no  & no  & no  \\
    A1 & prompt                 & yes & no  & no  \\
    A2 & deterministic          & no  & yes & no  \\
    A3 & supervisor             & no  & no  & yes \\
    A4 & prompt + deterministic & yes & yes & no  \\
    A5 & prompt + supervisor    & yes & no  & yes \\
    A6 & deterministic + supervisor & no & yes & yes \\
    A7 & all three              & yes & yes & yes \\
    \bottomrule
  \end{tabular}
\end{table}

\end{document}
